\begin{statement}{D.54}{Proved}{The local section theorem for the full geometrically prosolvable group}
Let \(k/\mathbb Q_p\) be finite and let \(U/k\) be a smooth
geometrically connected hyperbolic curve.  Every continuous section of
\[
  \pi_1(U)/\ker\bigl(\pi_1(U_{\bar k})\to
              \pi_1(U_{\bar k})^{\mathrm{sol}}\bigr)
                  \longrightarrow G_k
\]
is induced by a \(k\)-rational point of \(U\), or belongs to the
packet of a \(k\)-rational cusp.  If \(U\) is proper, only the first
alternative occurs.  The quotient is the full prosolvable inverse
limit, with no bound on derived length.  The proof uses the checked
analytic constructions of the supplied local preprint and proves
that they retain solvable monodromy; it does not assume its main
section theorem or a prosolvable variant of that theorem.
\end{statement}
\begin{proof}
We first prove the assertion for a proper hyperbolic curve
$X/k$. The affine case is treated in Step 7 by its actual
finite solvable section neighborhoods.
\paragraph{Step 1: exact inputs and fixed two-prime set.}

Fix once and for all an odd prime $\ell\ne p$ and put
$\Sigma=\{p,\ell\}$.  This same set is used on every curve and in every
comparison below.  The primary external statements needed are:

\begin{enumerate}
\item Mochizuki--Tsujimura, \emph{Resolution of Nonsingularities,
Point-theoreticity, and Metric-admissibility for p-adic Hyperbolic
Curves}, 30 March 2024, Theorem A and Proposition 2.4(iv):
geometrically pro-$\Sigma$ finite covers provide stable-model domination \cite{MT2024}.
Theorem F(i), equivalently Proposition 2.4(vii), of \cite{MT2024} gives a unique
section-fixed point in the universal pro-$\Sigma$ pro-Berkovich space
of a proper hyperbolic curve when the inertia and cyclotomic hypotheses
there hold.  Both hold for $H=G_k$: tame inertia surjects onto
$\mathbb Z_\ell$, and the closure of the powers of the residue-field
cardinality is open in $\mathbb Z_\ell^\times$.
\item Amini--Baker--Brugall\'e--Rabinoff, \emph{Lifting harmonic
morphisms I}, arXiv:1303.4812v3, Corollary 6.21 and Theorem 7.4(2):
the automorphisms of a lifted cover are precisely the stabilizer of
its annular gluing data \cite{DE_ABBR2015}.  Theorem 7.7 permits subdivision of loops.
We construct invariant gluing data below; existence of an arbitrary
lift alone would not establish the required Galois group.
\item Hoshi, \emph{Conditional results on the birational section
conjecture over small number fields}, RIMS-1742, Proposition 1.7:
over $k$, every geometrically pro-$\mathcal C$ birational section is
geometric if the full formation $\mathcal C$ contains the $p$-groups.
We apply this to finite solvable groups \cite{DE_Hoshi2012}.  In particular the conclusion
concerns the whole section image in that quotient, not only its
further pro-$p$ image.
\end{enumerate}

The analytic calculations of Sections 3--6 of the supplied local
preprint \cite{A_OpenAILocal} are used with their stated domains: roots exist on the
\emph{entire} inverse collar, and the transfer is of an $H$-valued
point, including its residue-field embedding, as in \cite[Fact 5.1.3]{DE_Temkin2010}.  Steps 3--5 below
check the additional monodromy assertions and recall the relevant
analytic calculations.  No prosolvable version of the supplied
preprint is being quoted.

\paragraph{Step 2: solvable covers and a compatible fixed point.}

Write $\Delta=\pi_1(X_{\bar k})^{\mathrm{sol}}$ and
$\Pi=\pi_1(X)/\ker(\pi_1(X_{\bar k})\to\Delta)$, and fix an actual
section $s:G_k\to\Pi$.  Since a finitely generated profinite group has
only finitely many open subgroups of each prescribed index, the
characteristic open subgroups $V\subset\Delta$ are cofinal.  The open
subgroup
\[
                 H_V=V\,s(G_k)\subset\Pi
\]
defines a connected finite \'etale cover $X_V\to X$ over $k$.
It is geometrically connected, and its geometric cover is Galois
with group $\Delta/V$.

The geometrically prosolvable fundamental group of $X_V$ is exactly
$H_V$.  Here is the needed finite-group argument.  A finite solvable
Galois cover of $X_{V,\bar k}$ has a normal closure over $X_{\bar k}$.
Over the Galois cover $X_{V,\bar k}$ this closure is a compositum of
finitely many conjugate solvable extensions.  Its Galois group over
$X_{V,\bar k}$ is a subgroup of the product of their solvable groups;
its quotient over $X_{\bar k}$ is $\Delta/V$, also solvable.  The total
group is therefore solvable.  This proves the asserted identification
by the classification of finite covers.  The same argument applies to
finite function-field extensions, and will also be used below.

Burnside's two-prime theorem implies that every finite $\Sigma$-group
is solvable.  Thus the section of $H_V$ projects to an actual section
$s_V^\Sigma$ of the geometrically pro-$\Sigma$ group of $X_V$.
By the first external input, it has a unique fixed point $z_V$ in
the universal pro-$\Sigma$ pro-Berkovich space of $X_V$.

These points are compatible, for a precise reason that does not
compare unrelated quotients.  If $V'\subset V$, pullback of a
geometric finite $\Sigma$-Galois cover of $X_V$ to $X_{V'}$ again has
$\Sigma$ monodromy.  With the common geometric base point, this
constructs the natural map from the universal pro-$\Sigma$ tower of
$X_{V'}$ to that of $X_V$.  The map is equivariant for the two sections,
since both sections are induced by the same homomorphism $s$.
The image of $z_{V'}$ is consequently fixed by $s_V^\Sigma$ in the
\emph{target} tower.  Its uniqueness there gives
\[
                         z_{V'}\longmapsto z_V.
\]
In particular their projections
$x_V\in (X_V\otimes_k\mathbb C_p)^{\mathrm{an}}$ are compatible.
The geometric covers $X_{V,\bar k}$ are cofinal in the universal
prosolvable geometric tower.  Hence $(x_V)_V$ is a point of that
whole pro-Berkovich tower, fixed by $s(G_k)$.

If its projection $x_X$ is type 1, every $x_V$ is type 1, since a
finite morphism has type-1 fibres over type-1 points.  The descent
action of $G_k$ fixes $x_V$, so $x_V\in X_V(k)$, using
$\mathbb C_p^{G_k}=k$, as proved in \cite[Theorem 1, p.~176]{DE_Tate1967}.  The compatible rational points in the
geometrically Galois covers show directly that $s(G_k)$ stabilizes
one lift of a point of $X(k)$.  Horizontal inertia in an \'etale
tower is trivial, so this stabilizer is precisely its point-section
subgroup.

If $x_X$ is not type 1, its restriction to the rational function
field is a real-valued rank-one valuation extending the $p$-adic
valuation: on a smooth analytic curve over $\mathbb C_p$, the zeros
and poles of a nonzero rational function are type 1.  The compatible
points give such a valuation on the universal solvable \'etale
field, preserved by $s(G_k)$.  The remaining argument excludes this
second alternative.

\paragraph{Step 3: cyclic roots on stable edges, with invariant gluing.}

Let $C/\mathbb C_p$ be proper of genus at least two, with split stable
skeleton, and let $e$ be one of its node edges.  We claim that there
is a finite \'etale cover of $C$ with solvable Galois closure on whose
entire inverse node annulus every annular coordinate has an analytic
$\ell$-th root.  In fact the edge construction below uses towers
of cyclic degree-$\ell$ covers.

Suppose first that $e$ lies on a simple circuit, oriented cyclically.
At each circuit vertex let $P$ be the outgoing mark and $Q$ the
incoming mark on its smooth residue curve $D$.  Divisibility of
$\operatorname{Pic}^0(D)$ gives a divisor $E$ such that
$P-Q+\ell E$ is principal.  The equation
\[
                          y^\ell=f,
       \qquad \operatorname{div}(f)=P-Q+\ell E
\]
defines a connected cyclic residue cover, totally ramified just at
$P,Q$.  Equip it with the generator $\tau:y\mapsto\zeta_\ell y$.
At every noncircuit vertex take $\ell$ identity copies, cyclically
permuted by $\tau$.  Circuit edges have one lift of expansion factor
$\ell$; every other edge has $\ell$ lifts of factor one.  Match the
unramified marked fibres by $C_\ell$-equivariant bijections.  Thus all
the data carry one coherent $C_\ell$ action.  The covering is
connected, and the local Riemann--Hurwitz equality
\[
  2g(D')-2=\ell(2g(D)-2)+2(\ell-1)
\]
shows that its augmented graph ramification is zero.  There are no
punctures, so a tame lift is \'etale.

The lift can be chosen $C_\ell$-equivariant.  On an orbit of
unramified edges, choose one gluing isomorphism and transport it
under $\tau$; the orbit has length $\ell$, so the last compatibility
is automatic.  For a circuit edge, its two branch coordinates can
be taken to have actions
\[
              u\longmapsto\zeta_\ell u,
       \qquad v\longmapsto\zeta_\ell^{-1}v.
\]
The inverse characters follow from the respective orders congruent
to $1$ and $-1$ modulo $\ell$ of $f$; multiplying a local Kummer
root by an integral power of the base parameter does not change
its deck character.  On the annular overlap the gluing has $uv=c$;
it commutes with these actions.  More intrinsically, the local
degree-$\ell$ annular deck group is $\mu_\ell$, and the two
restrictions are the same element when expressed in the common
oriented annular coordinate.  Thus the gluing data are fixed by
$\tau$.  Corollary 6.21 and Theorem 7.4(2) of the second external
input now lift the action.  Since the connected cover has degree
$\ell$, it is a $C_\ell$-Galois cover.  Subdividing loops inserts
the map $u\mapsto u^\ell$ and preserves these inverse characters.
The expansion factor on every annulus above $e$ is $\ell$.

If $e$ is separating, each side of the cut stable graph contains
a circuit or a positive-genus vertex.  A genus-zero tree side with
$N$ vertices would have total valence $2N-1$, whereas stability
requires at least $3N$.  On each side choose a connected cyclic
degree-$\ell$ unramified cover, using respectively a nonzero
homomorphism of the graph fundamental group to $C_\ell$, or a
nonzero $\ell$-torsion line bundle on the positive-genus residue
component.  Complete the residue and graph data by identity copies
and equivariant matching.  The $\ell$ bridge preimages can be
matched equivariantly between the two connected sides.  All gluing
edges are then free $C_\ell$-orbits, so the preceding transport
argument gives a connected cyclic \'etale cover of $C$.
Every bridge preimage is now nonseparating: the other $\ell-1$
bridges still join the two connected sides.  Apply the previous
construction to the finitely many bridge preimages and take a
connected component of the fibre product.  Its Galois closure has
an $\ell$-group as group, by the finite-group argument in Step 2.

The root assertion holds on whole annuli.  An analytic annular unit
has a factorization $aS^d(1+h)$ with $|h|<1$ everywhere on the open
annulus and a single dominant Laurent monomial.  Expansion factor
$\ell$ makes $d$ divisible by $\ell$; constants have roots in
$\mathbb C_p$ and $(1+h)^{1/\ell}$ converges on every compact
subannulus.  These roots agree on overlaps.  The inverse-skeleton
property of the tame lift identifies the full inverse node annulus
with the indicated node annuli, including their attached open
balls.  This proves the claim for one edge, and finite fibre
products give it for finitely many edges simultaneously.

For small collars around arbitrary algebraic points, perform the
blowups used in the supplied preprint to insert target node annuli.
Apply stable-model domination with our fixed $\Sigma=\{p,\ell\}$,
so its geometric covering group is solvable.  Choose an irrational
level strictly inside each target node annulus.  A source model
vertex gives a rational value to the target parameter; a smooth
disk tube also gives a rational constant value if it maps into that
target annulus.  Thus the full inverse irrational level lies in the
finitely many source node annuli.  Compactness gives an open collar
about the level whose full inverse lies in those annuli.  Apply the
edge construction to these finitely many annuli.  It follows that
the collar-root lemma of the supplied preprint holds with
\emph{solvable Galois closure} throughout.

\paragraph{Step 4: equivariant patching and arbitrary-order auxiliary roots.}

Let $C\to X_{\bar k}$ be a finite Galois cover with specified finite
solvable group $G$.  At every branch point $b$, choose a sufficiently
small disk with parameter $t_b$ such that each inverse disk has the
normal form $t_b=z^{e}$.  This uses characteristic zero and analytic
roots of units after shrinking; it does not require $e$ prime to $p$.
Suppose a finite cover $Y\to X_{\mathbb C_p}$ supplies an analytic
$m_b$-th root of $t_b$ on the entire inverse collar, with each
$e\mid m_b$.  Outside the inner disks the pullback of $C$ is a
$G$-torsor.  On each inverse disk replace it by the trivial
$G$-torsor.  On the overlap the equations split into all linear
factors, so the exterior torsor is trivial there.  A choice of one
sheet on each connected overlap identifies it \emph{equivariantly}
with $G$; the other sheets are its $G$-translates.  Glue through
these $G$-equivariant trivializations.

The result is a finite \'etale $G$-torsor on $Y^{\mathrm{an}}$.
There is no compatibility between distinct branch disks because
they are disjoint.  Proper analytic GAGA \cite[Theorem A.1(ii)]{DE_Poineau2010} algebraizes both the
cover and its $G$ action; \cite[Fact 5.1.5]{DE_Temkin2010} detects etaleness after analytification.  A connected component $V$ is a Galois
cover of $Y$ with group a subgroup of $G$, hence solvable.  If
$Y\to X$ has solvable Galois closure, so does $V\to X$, by
Step 2.  Descent of a finite \'etale cover along the algebraically
closed extension $\bar k\subset\mathbb C_p$ preserves the action
by full faithfulness.  No descent of the exterior analytic map is
needed.

For arbitrary positive $m$, the auxiliary construction of Section 5
of the supplied preprint also has solvable Galois closure.  Its
first curve is
\[
                  Z:\quad y^\ell=\frac{T(T-1)}{T-\alpha}.
\]
This is a cyclic degree-$\ell$ cover, with horizontal indices
$\ell$ at $0,\alpha,1,\infty$.  Its stable skeleton consists of
two rational vertices joined by $\ell$ parallel edges, each of
length $a=v_p(\alpha)$, and $g(Z)=\ell-1$.  Let $J$ be its
Jacobian and $i:Z\to J$ an Abel--Jacobi map.  The cover
\[
                       Z'=Z\times_{i,J,[m]}J
\]
is a torsor under the abelian finite group $J[m]$ over $\bar k$.
Every connected component has abelian Galois group over $Z$.
Its normal closure over the $T$-line therefore has solvable group,
an extension of an abelian subgroup of a product by $C_\ell$.
Its horizontal ramification indices still divide $\ell$.

For completeness, the roots supplied by this cover are the required
global analytic roots on all inverse sheets.  The totally degenerate
Jacobian has toric uniformization $\mathbb T^{\mathrm{an}}/\Lambda$ \cite[Sections 4.1--4.2 and 4.5.2]{DE_BakerRabinoff2015}.
The Abel--Jacobi map on each of the $\ell$ open annuli lifts to
$\mathbb T^{\mathrm{an}}$.  The cycle $e_j-e_{j'}$ gives a character
whose pullback has slope one by \cite[Proposition 6.1 and Corollary 6.6]{DE_BakerRabinoff2015}, hence is
\[
                 u_j=c_j T(1+h_j),\qquad
        |h_j|\le p^{-\min\{v_p(T),a-v_p(T)\}}.
\]
On $Z'$ the character has an analytic $m$-th root.  Indeed local
torus lifts of the division point have $m$-th power equal to the
original lift times a period.  Decompose by its coset in
$\Lambda/\Lambda^m$, choose one period representative, and adjust
local lifts to that representative.  Their differences on overlaps
are $m$-torsion periods, hence trivial because $\Lambda$ is torsion
free.  The adjusted lifts therefore glue on every clopen piece.
If
\[
        D>v_p(m)+\frac1{p-1},\qquad a>2D,
\]
then $\exp(m^{-1}\log(1+h_j))$ converges on
$D<v_p(T)<a-D$.  Dividing the character root by this function and
an $m$-th root of $c_j$ produces a root of $T$ on the whole inverse
middle collar.  Pullback by $T=t_b/\beta_b$ and finite composita
retain both this property and solvability of the Galois closure.

We also use the following horizontal ramification fact.  Over the
algebraically closed characteristic-zero residue field of a
horizontal valuation, the local extensions of the completed
function field have the form $\bar k((t^{1/e}))$.  Consequently
indices dividing $\ell$ remain so under composita, conjugates,
normal closures, normalized pullback, and composition on top with
an \'etale cover.  Thus the normal closures just used do not lose
the horizontal index bound.

\paragraph{Step 5: completion contains every solvable birational extension.}

Put $K=k(X)$ and $K_0=\bar k(X)$.  Inside a fixed $K^{\mathrm{sep}}$,
let $B$ be the maximal solvable extension of $K_0$, and let $L\subset B$
be the universal geometrically solvable \'etale field of $X$.
Both are normal over $K$.  Fix any prolongation of a real rank-one
valuation extending $v_p$, embed their completions in
$E=\widehat{K^{\mathrm{sep}}}$, and put $H=\widehat L\subset E$.
In particular $\mathbb C_p\subset H$.  We prove
\[
                              B\subset H.                 \tag{*}
\]

First let $F/K_0$ be finite Galois with solvable group and with every
horizontal index dividing $\ell$.  Its branch disks can be chosen
to avoid the analytic point defined by $K_0\hookrightarrow H$:
this point projects to the generic point of $X_{\bar k}$, and hence
is distinct from each algebraic branch point.  Step 3 gives an
\'etale cover $Y\to X$ with solvable Galois closure and all the
necessary collar roots.  Equivariant patching in Step 4 gives
an \'etale cover $V\to X$ with solvable Galois closure and an
exterior analytic map $V\to C$, where $C$ has function field $F$.
Its function field embeds in $L$, and hence its $H$-valued generic
point transfers through the exterior analytic map to an $H$-valued
generic point of $C$.  On completed residue fields this is a
$K_0$-embedding $F\hookrightarrow H$.  Normality of $F/K_0$ makes
its image equal to the chosen field $F\subset E$.  This is an
inclusion of the field, not only the existence of an abstract
isomorphic copy.

Let $L_\ell^{\mathrm{sol}}$ be the compositum of all these finite
solvable Galois extensions with indices dividing $\ell$.  The first
stage proves $L_\ell^{\mathrm{sol}}\subset H$.

Now take \emph{any} finite solvable Galois extension $F/K_0$.
At a branch point $b$ let $m_b$ be the least common multiple of its
indices.  The auxiliary covers of Step 4, pulled back by the
local rational parameters and combined by a compositum, give
$W\to X_{\bar k}$ with solvable Galois closure, horizontal indices
dividing $\ell$, and all analytic $m_b$-th roots on the full inverse
collars.  Equivariant patching produces $V\to W$ \'etale with
solvable group on each connected component.  Thus $V\to X$ has
solvable Galois closure and horizontal indices dividing $\ell$.
Its field embeds in $L_\ell^{\mathrm{sol}}\subset H$.  The same
$H$-valued point transfer and normality again give $F\subset H$.
Every element of $B$ lies in one of these finite solvable Galois
extensions.  This proves (*), with no bound on the ramification
orders or derived lengths of $F$.

Restriction of decomposition groups is consequently an isomorphism
\[
                  D_w(B/K)\xrightarrow{\sim}D_w(L/K).       \tag{**}
\]
Surjectivity follows by lifting an automorphism and adjusting it by
$\operatorname{Gal}(B/L)$, which acts transitively on prolongations;
the infinite case follows from finite-level transitivity and
compactness.  An element of the kernel is an isometry after
normalization $w(p)=1$ and extends to $\widehat B$.  It fixes
$H=\widehat L$ pointwise.  But (*) implies
$B\subset H\subset\widehat B$, so it fixes $B$.  This proves
injectivity; compactness makes the inverse continuous.

\paragraph{Step 6: excluding a section in a rank-one decomposition group.}

If $s(G_k)$ fixes a rank-one valuation in the \'etale tower,
(**) gives a continuous lift
\[
             \widehat s:G_k\longrightarrow\operatorname{Gal}(B/K)
\]
into its birational decomposition group.  By Hoshi's Proposition 1.7
the whole image also lies in a horizontal decomposition group at
some $a\in X(k)$.  Let
$M=B^{\widehat s(G_k)}$.  The restrictions $v_1,v_2$ of the vertical
and horizontal valuations to $M$ have rank one, and
\[
                         v_1(p)=1,\qquad v_2|_k=0.
\]
Rank one is preserved here because all the field extensions from
$K$ are algebraic, so their value-group quotients are torsion.

The required Hensel argument does not presume that $M$ is henselian.
The field $B$ is quadratically closed.  In fact for any $b\in B$,
choose a finite solvable Galois extension $F/K_0$ containing $b$;
adjoining square roots of all its conjugates is an abelian extension
of $F$ whose normal closure over $K_0$ is solvable.  Thus $\sqrt b\in B$.
Every quadratic extension of $M$ is therefore contained in $B$.
Since all of $\operatorname{Gal}(B/M)=\widehat s(G_k)$ stabilizes
each selected valuation, every such quadratic extension has a
unique prolongation of $v_i$.  Equivalently these valuations are
2-henselian.  More concretely, if $u\in M$ satisfies
\[
                             v_i(u-1)>2v_i(2),
\]
Hensel's lemma gives a square root in the henselization of $M$ at
$v_i$.  It also belongs to $B$.  On the quadratic extension
$M(\sqrt u)/M$, the full Galois group is its decomposition group,
so its intersection with that henselization is $M$.  Hence
$\sqrt u\in M$.  This proves exactly the square-root assertion
used next, including $p=2$ and nondiscrete valuations.

Choose a nonsquare $d\in k^\times$.  It remains nonsquare in $M$,
since the projection $\widehat s(G_k)\to G_k$ is onto.  Let $t\in K$
be a uniformizer at $a$, choose an integer $N>v_1(t)$ and set
$z=p^N/t$.  Then $v_1(z)>0$ and $v_2(z)<0$.  Put
\[
                    q_r=d+(1-d)\frac{z^r}{1+z^r}.
\]
The denominator is nonzero, and
\[
\begin{aligned}
 v_1(q_r/d-1)&=v_1((1-d)/d)+r v_1(z)\longrightarrow+\infty,\\
 v_2(q_r-1)&=v_2(d-1)-r v_2(z)\longrightarrow+\infty.
\end{aligned}
\]
For large $r$, the preceding square-root assertion gives
$\sqrt{q_r/d}$ and $\sqrt{q_r}$ in $M$.  Their quotient gives
$\sqrt d\in M$, a contradiction.  Thus no section fixes a
rank-one valuation in $L$.  Combined with Step 2, this proves
the proper-curve assertion.

\paragraph{Step 7: affine passage.}

For an affine hyperbolic curve $U$, the following elementary
cofinality fact supplies proper hyperbolic compactifications inside
its solvable tower.  In the free pro-$r$ geometric fundamental
group, each peripheral element has infinite order, for any fixed
prime $r$: with at least two punctures it has nonzero primitive
abelian class, and with one puncture it is the nontrivial product
of surface commutators.  Its orders in finite $r$-group quotients
are therefore unbounded.  Choose a quotient in which all the
finitely many peripheral orders exceed a common $M$, and take a
characteristic refinement.  If its degree is $d$ and the boundary
has $n$ points, the compactification has
\[
     2g'-2=d\left(2g-2+n-\sum_{j=1}^{n}\frac1{e_j}\right).
\]
Since $2g-2+n>0$, sufficiently large $e_j$ make $g'\ge2$.
The same construction over each finite solvable section-neighborhood,
followed by its solvable normal closure over $U$, gives cofinally
many such neighborhoods.  The section-induced arithmetic twists
are the covers defined by $V s(G_k)$; hence they are over $k$ and
the section lifts to them.

The section on each such affine cover induces one on the
geometrically prosolvable group of its proper compactification.
The proper-curve assertion gives a $k$-rational point there.
The compact spaces of these rational points have nonempty inverse
limit: any finite collection is dominated by a further section
neighborhood, whose rational point supplies compatible images.
The resulting point in the compactified solvable tower is
section-fixed.  Its stabilizer is the ordinary point decomposition
group if the base point lies in $U$, and the cusp decomposition
group otherwise.  This gives exactly point/cusp origin.

\end{proof}
